% Proposed isolated A.91. Actual marked modules, not generic parameters.
\statement{A.91}{Proved}{The exact abelian scope of the marked cube reversal}
Retain the actual \(\widehat{GT}_0\) data of A.87 and put
\[
 T_d=v\delta^d,\qquad \Xi_d=T_d^3,\qquad
 \sigma_d=\operatorname{Ad}(\delta^{-d})\Phi_\alpha|_F.
 \tag{A.91.1}
\]
For \(i=x,y,z\), let
\[
 K_i=\overline{\langle\!\langle i\rangle\!\rangle}_F,\qquad
 M_i=K_i/\overline{[K_i,K_i]},\qquad
 M=M_x\oplus M_y\oplus M_z.
 \tag{A.91.2}
\]
Then \(\sigma_d(\Xi_d)=\Xi_d^{-1}\), but the induced action
of \(\sigma_d\) on the entire \(M_x\) is \(-\mathrm{id}\).
Thus this reversal equality is tautological on every
invariant abelian coefficient quotient of \(M_x\).

More precisely, let \(\mathsf t\) be conjugation by \(T_d\)
on the three-cover sum M. The actual \(\mathsf t,\tau\)
actions satisfy
\[
 \mathsf t^3=1,\qquad \tau^2=1,\qquad
 \tau\mathsf t\tau=\mathsf t^{-1}.
 \tag{A.91.3}
\]
If \(\jmath_x:M_x\to M\) is the summand inclusion, the map
\[
 a\longmapsto
 \jmath_x(a)+\mathsf t\jmath_x(a)+\mathsf t^2\jmath_x(a)
 \tag{A.91.4}
\]
is a continuous isomorphism onto
\[
             \{w\in M:\mathsf t(w)=w,\ \tau(w)=-w\}.
 \tag{A.91.5}
\]
In particular these actual dihedral covariance identities
impose no further restriction on a single cover coordinate.
No assertion is made that every such vector comes from one
common full word or a canonical hexagon defect.

\begin{proof}
\emph{1. The literal reflection.}\par
A.87 proves \(\tau(T_d)=T_d^{-1}\), \(\Xi_d\in F'\),
and \(\tau(\Xi_d)=\Xi_d^{-1}\).
Because \(\Xi_d\) commutes with \(T_d\), the product
\(T_d^{-1}\tau\) also reverses it. In the actual semidirect
product its ordered expression is
\[
                  T_d^{-1}\tau
                    =\delta^{-d}vv\alpha=\delta^{-d}\alpha.
 \tag{A.91.6}
\]
This proves the asserted \(\sigma_d\) equation.
With \(m=(d-1)/2\), the equality
\(\delta^{-d}=x^m\delta^{-1}\) gives
\[
 \sigma_d=\operatorname{Inn}(x^m)\sigma_1,\qquad
 \sigma_1=\operatorname{Ad}(\delta^{-1})\Phi_\alpha|_F.
 \tag{A.91.7}
\]
The original-marking formula in A.87 is
\[
 \sigma_1(x)=x^{-1},\qquad
 \sigma_1(y)=j_\delta(f)^{-1}xyj_\delta(f),\qquad
 j_\delta=\operatorname{Ad}(\delta^{-1}).
 \tag{A.91.8}
\]
Thus both \(\sigma_1\) and \(\sigma_d\) preserve \(K_x\)
and induce the identity on its procyclic y quotient.
The actual d is odd. A.45 further gives m even, so one
also has the equality of existing full automorphisms
\[
 \sigma_d
   =\operatorname{Ad}(x^{m/2})\sigma_1
        \operatorname{Ad}(x^{-m/2}).
 \tag{A.91.9}
\]

\emph{2. The full integral action.}\par
The compact abelian group \(M_x\) is topologically generated
by the classes of \(y^\lambda x y^{-\lambda}\),
\(\lambda\in\widehat{\mathbb Z}\). Indeed K is the closed
normal closure of x, \(F/K_x\) has its literal y section,
and the K factor of any conjugator acts trivially on the
abelianization. Since \(\sigma_d(y^\lambda)\) has the same
y head and \(\sigma_d(x)=x^{-1}\), one obtains
\[
 [\sigma_d(y^\lambda x y^{-\lambda})]
       =-[y^\lambda x y^{-\lambda}].
 \tag{A.91.10}
\]
Continuity proves \(\sigma_d=-\mathrm{id}\) on the whole
\(M_x\). No free-module presentation is needed.

Accordingly the equality on \(\Xi_d\) becomes
\(-[\Xi_d]=-[\Xi_d]\). Since
\(\mathcal H_d=u\Xi_d^{-1}u^{-1}\) and \(u\in K_x\),
its integral cover class is \(-[\Xi_d]\), and its
mod-two class \(R_d\) equals \([\Xi_d]\).
The reversal therefore gives only \(R_d=R_d\).
It does not give A.89's distinct weighted condition
\(\overline{R_d}=T^dR_d\).

\emph{3. The complete three-cover action.}\par
Conjugation by \(\delta\) fixes \(K_x\) and exchanges
\(K_y,K_z\). Conjugation by v exchanges \(K_x,K_y\)
and fixes \(K_z\), since v is the actual inner-conjugated
swap. Thus, for odd d, \(T_d=v\delta^d\) cycles the
three normal subgroups in the order x,y,z.
The actual peripheral involution \(\alpha\) preserves
each \(K_i\), so \(\tau=v\alpha\) exchanges \(K_x,K_y\)
and fixes \(K_z\).

These induce continuous operators on M. Their cube is
conjugation by \(\Xi_d\). Since
\(\Xi_d\in F'\subseteq K_i\) for every i, that inner
conjugation is trivial on every \(M_i\).
The other two relations in (A.91.3) follow from the actual
involution and reversal identities before abelianization.
The stabilizer reflection \(\mathsf t^{-1}\tau\) on
\(M_x\) is exactly (A.91.6), so acts as \(-\mathrm{id}\).
Hence
\[
                 \tau\jmath_x(a)=-\mathsf t\jmath_x(a).
 \tag{A.91.11}
\]
Using \(\tau\mathsf t^j=\mathsf t^{-j}\tau\)
shows that \(\tau\) negates the sum in (A.91.4);
\(\mathsf t\) fixes it. Conversely a \(\mathsf t\)-fixed
vector is uniquely determined by its x coordinate because
\(\mathsf t\) cycles the three summands. The x-coordinate
projection is the continuous inverse, proving (A.91.4)--(A.91.5).

The common-word homomorphism \(F'\to M\) sends the actual
\(\Xi_d\) to a vector in (A.91.5). The theorem does not
identify the image of that homomorphism with all of
(A.91.5), nor discard the full pentagon or canonical-word
constraints. For general \(\widehat{GT}_0\) data it adds
no missing odd-head identity to A.89; for full-GT data
the dyadic conclusion remains proved while A.90's
all-odd-head condition stays open. The nonlinear
\(\Xi_d/\Omega_\alpha\) equation remains separate.
\end{proof}
